\section{BiTempQA Benchmark}

\subsection{Problem Formulation}
Following bitemporal data modeling~\citep{snodgrass1999tsql2}, we define two temporal dimensions for memory-system evaluation.
\textit{Event time} ($t_e$) is when an event actually occurred in the real world.
\textit{Record time} ($t_r$) is when the memory system recorded the event.
A memory entry is represented as a tuple $(\textit{text}, t_e, t_r)$, and a memory store at logical time $T$ contains all entries with $t_r \leq T$. A question may require reasoning about $t_e$ only, $t_r$ only, or both simultaneously.

\subsection{Data Sources}
A core design principle of BiTempQA is grounding scenarios in real-world temporal data rather than purely synthetic generation. Our dataset combines two data sources:

\paragraph{ChronoQA~\citep{chronoqa2025}.}
We extract 200 bitemporal candidates from ChronoQA, a Chinese temporal QA dataset of news events with natural event-time/record-time separation (event date vs.\ publication date).
Each candidate includes golden text chunks, event dates, and publication dates, which we map directly to the ($t_e$, $t_r$) dimensions.
We select candidates with gap $\geq 1$ day between event and publication, yielding a balanced distribution: 100 with 1--6 day gaps, 40 with 7--29 day gaps, 30 with 30--364 day gaps, and 30 with 365+ day gaps.

\paragraph{Curated real-world events.}
We manually construct 15 scenarios from well-documented public events across diverse domains: CEO transitions (Microsoft, Apple, Alibaba), political leadership changes (Japanese/UK Prime Ministers), sports coaching changes (Real Madrid), corporate headquarters migration (Lenovo), scientific discoveries (Pluto reclassification, Tu Youyou Nobel Prize, Chang'e lunar missions), health events (COVID-19 timeline, vaccine development), and multi-source news events (MH370, SpaceX-NASA collaboration).
Each scenario includes 3--7 memory writes with verified dates from public records.

\paragraph{LLM augmentation.}
For scenario types not naturally covered by real data (contradictory information, future-dated information, entity identity resolution, temporal ambiguity), we retain 118 LLM-generated scenarios from an initial template-based generation round using Qwen2.5-72B-Instruct.
This yields a dataset that is 64.6\% real-source (215 scenarios) and 35.4\% LLM-generated (118 scenarios).

\subsection{Scenario Design}
We design 10 scenario types spanning 6 domains (corporate, academic, social, fictional, historical, and scientific):
(1) Entity Attribute Evolution: attributes change over time (e.g., job titles);
(2) Relationship Evolution: social/professional relationships change;
(3) Contradictory Information: different sources report conflicting facts;
(4) Late-Arriving Facts: events recorded well after they occurred ($|t_r - t_e|$ large);
(5) Future-Dated Information: events recorded before their stated occurrence ($t_r < t_e$);
(6) Entity Identity Resolution: same entity referenced differently across time;
(7) Knowledge Retraction: previously recorded facts are later revoked;
(8) Multi-Source Information: multiple independent sources report overlapping events;
(9) Gradual Accumulation: knowledge builds incrementally over time;
(10) Temporal Ambiguity: unclear or ambiguous temporal references.

Each scenario consists of 2--7 memory writes, each annotated with explicit ISO~8601 timestamps for $t_e$ and $t_r$. In most cases, scenarios are constructed so that $t_e \neq t_r$, creating natural bitemporal tension.

\subsection{Question Taxonomy}
Questions are organized into 3 types at 2 difficulty levels:

\textit{Level 1 (Easy).} Point-in-time retrieval (58.3\%) and temporal ordering (21.4\%).

\textit{Level 2 (Medium).} Multi-hop temporal reasoning (20.2\%).

All questions use abstractive (free-text) answer format. Each question is annotated with metadata indicating whether it requires event time reasoning, record time reasoning, version tracking, or knowledge retraction.

\subsection{Dataset Statistics}
BiTempQA v3 contains 415 QA pairs across 333 scenarios (215 real-source + 118 LLM-generated), split into train (290, 70\%), dev (41, 10\%), and test (84, 20\%).
All reported results use the held-out test set, which consists entirely of real-source scenarios (64 ChronoQA-sourced + 20 from curated events).
Table~\ref{tab:composition} shows the dual-timestamp composition of the test set.

\begin{table}[t]
\centering
\resizebox{\columnwidth}{!}{%
\begin{tabular}{lrl}
\toprule
\textbf{Temporal Requirement} & \textbf{Count (\%)} & \textbf{Example Question Type} \\
\midrule
Both timestamps required & 47 (56.0\%) & ``What was known about X as of date Y?'' \\
Event-time only & 37 (44.0\%) & ``When did event X occur?'' \\
Record-time only & 0 (0.0\%) & --- \\
\midrule
Version tracking needed & 0 (0.0\%) & --- \\
Knowledge retraction & 0 (0.0\%) & --- \\
\bottomrule
\end{tabular}%
}
\caption{Dual-timestamp composition of the BiTempQA v3 test set (84 questions). Over half of the questions require reasoning about both $t_e$ and $t_r$.}
\label{tab:composition}
\end{table}

Table~\ref{tab:scenario_dist} shows the distribution of scenario types across the full dataset.

\begin{table}[t]
\centering
\resizebox{\columnwidth}{!}{%
\begin{tabular}{lrrr}
\toprule
\textbf{Scenario Type} & \textbf{Real} & \textbf{LLM} & \textbf{Total} \\
\midrule
Entity Attribute Evol. & 114 & 0 & 114 \\
Multi-Source Information & 98 & 0 & 98 \\
Temporal Ambiguity & 35 & 0 & 35 \\
Late Arriving Facts & 2 & 14 & 16 \\
Knowledge Retraction & 1 & 12 & 13 \\
Contradictory Info. & 0 & 12 & 12 \\
Future-Dated Info. & 0 & 12 & 12 \\
Gradual Accumulation & 1 & 11 & 12 \\
Relationship Evol. & 1 & 10 & 11 \\
Entity Identity Resol. & 0 & 10 & 10 \\
\midrule
\textbf{Total} & \textbf{215} & \textbf{118} & \textbf{333} \\
\bottomrule
\end{tabular}%
}
\caption{Scenario type distribution by data source. Real-source scenarios concentrate on entity evolution, multi-source information, and temporal ambiguity; LLM-generated scenarios cover rare types with no natural real-world data.}
\label{tab:scenario_dist}
\end{table}

\subsection{Evaluation Pipeline}
We implement a unified Retrieve $\rightarrow$ Generate $\rightarrow$ Judge pipeline (Figure~\ref{fig:framework}):

\textit{Step 1: Retrieve.} The memory backend retrieves the top-$k$ context passages given the question and the scenario memory store.

\textit{Step 2: Generate.} DeepSeek-V3~\citep{deepseekv3} generates an answer conditioned on the retrieved context.

\textit{Step 3: Judge.} An independent LLM call scores the generated answer against the gold answer. Since all test questions are abstractive, DeepSeek-V3 in ``answer\_judge'' mode evaluates semantic correctness with lenient matching. Without the LLM judge, exact text matching is too strict for abstractive evaluation; with the judge, accuracy reflects semantically appropriate correctness.
